% Style based on the supplied Waseem_Generic_AI_ML_GenAI_AgenticAI.tex.
% Python replaces template tokens. Edit layout here, content in JSON/portfolio JS.
\documentclass[9pt,letterpaper]{extarticle}
\usepackage[left=0.68in,right=0.68in,top=0.45in,bottom=0.45in]{geometry}
\usepackage[T1]{fontenc}
\input{glyphtounicode}
\pdfgentounicode=1
\usepackage[utf8]{inputenc}
\usepackage{lmodern}
\usepackage{microtype}
\usepackage{pifont}
\usepackage{xcolor}
\usepackage{enumitem}
\usepackage{titlesec}
\usepackage[hidelinks]{hyperref}
\usepackage{tabularx,array,parskip,ragged2e,needspace}
\pagestyle{empty}
\setlength{\parindent}{0pt}
\setlength{\parskip}{1pt}
\definecolor{accent}{HTML}{155B8A}
\titleformat{\section}{\large\bfseries\color{accent}}{}{0pt}{}[\vspace{0.5pt}\color{accent}\titlerule]
\titlespacing*{\section}{0pt}{3pt}{2pt}
% Compact subsections beneath the main accent-rule sections.
\titleformat{\subsection}{\normalsize\bfseries\color{accent}}{}{0pt}{}
\titlespacing*{\subsection}{0pt}{3pt}{1pt}
\setlist[itemize]{leftmargin=1.15em,itemsep=0.8pt,topsep=1pt,parsep=0pt,partopsep=0pt}
\newcommand{\role}[4]{%
  \Needspace{4\baselineskip}\textbf{#1}\hfill\textbf{#2}\\[-1pt]
  \textit{#3}\hfill\textit{#4}\par}
\newcommand{\entry}[2]{\Needspace{4\baselineskip}\textbf{#1}\hfill\textit{#2}\par}
\newcommand{\skillrow}[2]{\textbf{#1:} & #2\\}
\hypersetup{pdftitle={Waseem Raza, PhD - AI/ML for Wireless Networks},pdfauthor={Waseem Raza, PhD}}
\begin{document}

\begin{center}
{\LARGE\bfseries Waseem Raza, PhD}{\normalsize\mdseries }\\[2pt]
{\large\bfseries Machine Learning Engineer \textbar{} AI/ML for Wireless \& 5G/6G Systems}\\[1pt]
{\small Wireless AI \textbar{} 5G/6G \textbar{} RAN Intelligence \textbar{} Network Automation}\\[3pt]
{\footnotesize \ding{41}\enspace \href{mailto:waseem.placements@gmail.com}{waseem.placements@gmail.com} \textbar{} \href{mailto:waseemraza844@gmail.com}{waseemraza844@gmail.com}\\[1pt]New Jersey, United States \quad \ding{37}\enspace \href{tel:+17328039703}{+1 732-803-9703}\\[1pt]\href{https://www.linkedin.com/in/waseem-raza-phd-9a737858/}{LinkedIn} \textbar{} \href{https://github.com/WaseemRaza844}{GitHub} \textbar{} \href{https://scholar.google.com/citations?user=d58Drr0AAAAJ\&hl=en}{Google Scholar} \textbar{} \href{https://waseemraza844.github.io/portfolio/learning.html}{Coursera Courses} \textbar{} \href{https://waseemraza844.github.io/portfolio/}{Portfolio}}
\end{center}
\vspace{-5pt}
\Needspace{3\baselineskip}
\section{Professional Profile}
AI/ML Engineer for Wireless Networks with a PhD in Electrical and Computer Engineering and 7+ years of combined industry and research experience in 5G/6G systems, network intelligence, propagation modeling, mobility, RAN optimization, anomaly detection, and resilient automation. Apply deep learning, transfer learning, domain adaptation, generative modeling, reinforcement learning, and statistical methods to network-assurance, NR-NTN, CSI compression, digital-twin, outage-management, and self-healing problems. Experienced with Python, PyTorch, TensorFlow, scikit-learn, SQL, large-scale network measurements, 3GPP-compliant simulation, and production-oriented ML validation.
\Needspace{2\baselineskip}
\subsection{Technical Competencies}
\Needspace{2\baselineskip}\noindent{\small\textbf{AI/ML for Wireless:} Transfer learning, domain adaptation, adversarial learning, generative modeling, reinforcement learning, anomaly detection, classification, and robust model evaluation}\par
\Needspace{2\baselineskip}\noindent{\small\textbf{Wireless Systems:} LTE/4G, 5G-NR, NR-NTN, mobility, RAN, propagation, cell selection, user timing, network measurements, and service assurance}\par
\Needspace{2\baselineskip}\noindent{\small\textbf{Network Intelligence:} Root-cause analysis, outage detection and compensation, self-healing, optimization, resilience, and intelligent automation}\par
\Needspace{2\baselineskip}\noindent{\small\textbf{Simulation \& Digital Twins:} 3GPP-compliant simulation, propagation modeling, synthetic data generation, digital twins, large-scale experimentation, and robustness testing}\par
\Needspace{2\baselineskip}\noindent{\small\textbf{Programming \& Data:} Python, SQL, PyTorch, TensorFlow, scikit-learn, XGBoost, Pandas, NumPy, visualization, and statistical analysis}\par
\Needspace{2\baselineskip}\noindent{\small\textbf{Engineering \& MLOps:} Git, Docker, CI/CD, reproducible experimentation, APIs, cloud analytics, model validation, and production-oriented workflows}\par
\Needspace{4\baselineskip}\textbf{Leadership and Professional Strengths}\par
\begin{itemize}
\item \textbf{Ownership and accountability:} Lead AI/ML initiatives through requirements analysis, data preparation, model development, validation, stakeholder review, and deployment planning.
\item \textbf{Technical leadership:} Define technical approaches, evaluate alternatives, coordinate experimentation, and communicate findings to guide projects from ambiguous requirements to measurable outcomes.
\item \textbf{Cross-functional collaboration:} Work with software engineers, network specialists, researchers, geologists, and public-safety stakeholders to translate domain requirements into practical data and AI solutions.
\item \textbf{Continuous learning:} Apply ongoing study in generative AI, agentic AI, RAG, LLM evaluation, and cloud-native platforms through hands-on projects and professional coursework.
\item \textbf{Reliability and commitment:} Maintain reproducible workflows, careful validation, and attention to detail throughout experimentation and delivery.
\item \textbf{Adaptability and problem solving:} Combine critical thinking, systematic experimentation, and data-driven analysis to address unfamiliar technologies and complex operational problems.
\end{itemize}
\Needspace{3\baselineskip}
\section{Professional Experience}
\role{Machine Learning Engineer}{02/2025 - Present}{AT\&T, contractor via Innovatix Inc.}{Bedminster, NJ}
\begin{itemize}\interlinepenalty=10000
\item Lead AI/ML development for E911 and network-assurance diagnostics, applying NLP, domain adaptation, and root-cause classification to operational telecom workflows.
\item Develop multimodal diagnostic models using work-order text, location information, network context, and operational signals for VoNR and 5G service assurance.
\item Design cross-domain validation and distribution-shift evaluation workflows to improve reliability of ML models across heterogeneous network environments.
\item Collaborate with network engineers, software teams, and public-safety stakeholders on AI-enabled automation for critical telecommunications operations.
\end{itemize}
\role{Advanced-Degrees Intern}{06/2024 - 08/2024}{AT\&T Labs}{Austin, TX}
\begin{itemize}\interlinepenalty=10000
\item Analyzed 500K+ LTE/5G-NR diagnostic logs and applied ML methods to detect anomalies, reducing diagnostic time by 20\%.
\item Investigated 5G NR-NTN cell-selection behavior and improved decision algorithms by 15\% through data-driven analysis and validation.
\item Conducted structured analysis of large operational datasets and communicated technical findings to engineering teams supporting next-generation wireless systems.
\end{itemize}
\role{Advanced Wireless Technology Development Intern}{09/2022 - 01/2023}{Futurewei Technologies}{Schaumburg, IL}
\begin{itemize}\interlinepenalty=10000
\item Improved autoencoder-based channel-state-information compression by 25\% using transfer learning and systematic experimentation on large wireless datasets.
\item Conducted 50+ experiments to validate model performance, robustness, and generalization under changing data conditions.
\item Collaborated with research and engineering stakeholders to translate modeling results into recommendations for next-generation communication systems.
\end{itemize}
\role{Advanced-Degrees Intern}{06/2022 - 08/2022}{AT\&T Labs}{Bedminster, NJ}
\begin{itemize}\interlinepenalty=10000
\item Designed ML and statistical methods for analyzing users' timing in 5G standalone networks, improving analytical precision by 20\%.
\item Analyzed large-scale user trace data with SQL and visualization tools and communicated findings to senior technical leadership.
\item Converted raw operational traces into interpretable performance indicators to support diagnosis and network optimization.
\end{itemize}
\role{Graduate Research Assistant}{09/2019 - 07/2025}{AI4Networks, University of Oklahoma}{Tulsa, OK}
\begin{itemize}\interlinepenalty=10000
\item Designed reusable Python simulation and data-generation platforms for large-scale experimentation, including a 3GPP-compliant network simulator and an AI-enabled 5G testbed.
\item Built transfer-learning, adversarial-learning, GAN, autoencoder, and reinforcement-learning pipelines for sparse-data and cross-domain modeling.
\item Developed robustness evaluation frameworks for models operating under distribution shifts, missing data, and limited labeled samples; led 100+ experiments from problem formulation through validation.
\item Developed scalable self-healing and outage-management frameworks that improved modeled network resilience by 30\%.
\item Published peer-reviewed research, presented technical results internationally, and mentored junior researchers in ML implementation, experimental design, and technical writing.
\end{itemize}
\Needspace{3\baselineskip}
\section{Education}
\entry{PhD, Electrical and Computer Engineering}{09/2019 - 07/2025}
University of Oklahoma, Norman, Oklahoma, USA\par
\entry{MSc, Telecommunication Engineering}{08/2014 - 08/2016}
University of Engineering and Technology, Taxila, Pakistan\par
\entry{BSc, Telecommunication Engineering}{11/2010 - 06/2014}
University of Engineering and Technology, Taxila, Pakistan\par
\Needspace{3\baselineskip}
\section{Selected Projects \& Research}
\Needspace{2\baselineskip}
\subsection{Projects \& Applied Research}
\entry{SyntheticNET: 3GPP-Compliant Network Simulation}{Research Project \textbar{} 2020-2024}
Developed a Python-based, 3GPP-aligned simulation environment for realistic mobility, propagation, and cellular-network experimentation when measured data is sparse or costly. Supported controlled data generation and repeatable studies in propagation modeling, network optimization, digital twins, and resilient wireless AI.\par
{\small\textit{Tools \& Methods:} Python, 3GPP, Wireless Simulation, Mobility Modeling, Propagation Modeling, Synthetic Data, Experiment Design, Digital Twins, Network Analytics\par}
\vspace{1pt}
\entry{TurboRAN Experimental Platform}{Research Infrastructure \textbar{} 2020-2024}
Built and used a multi-band, multi-tier experimental platform for evaluating AI-enabled radio-access-network methods under controlled network conditions. Supported repeatable measurement collection, testbed experimentation, RAN optimization studies, and validation of emerging wireless-network algorithms.\par
{\small\textit{Tools \& Methods:} 5G Testbeds, RAN, Network Measurement, Experiment Design, Wireless Prototyping, Multi-Band Systems, Performance Evaluation, RAN Optimization\par}
\vspace{1pt}
\entry{Propagation Models \& Digital Twins}{Research Project \textbar{} 2024}
Developed robust data-driven radio-propagation models for wireless digital twins using domain-informed generative learning and scarce measured data. Combined 3GPP-informed pretraining, GAN-based augmentation, domain adaptation, and statistical validation to improve resilience under distribution shift.\par
{\small\textit{Tools \& Methods:} Propagation Modeling, Digital Twins, GANs, Domain Adaptation, 3GPP Models, Synthetic Data, Statistical Validation, Wireless AI, Distribution Shift\par}
\vspace{1pt}
\entry{AI for Telecom Network \& Security Optimization}{Certification Project \textbar{} 2026}
Applied AI concepts to telecommunications network optimization, security analysis, and customer-experience improvement through a structured specialization project. Framed telecom use cases around operational data, AI-assisted analysis, optimization opportunities, and practical service-improvement recommendations.\par
{\small\textit{Tools \& Methods:} AI for Telecommunications, Network Optimization, Telecom Security, Customer Experience, Python, TensorFlow, Operational Analytics, Service Improvement\par}
\vspace{1pt}
\entry{4G Cellular Network Fundamentals Lab}{Certification Project \textbar{} 2026}
Consolidated LTE/4G architecture, radio-access and core-network roles, mobility, and service-flow concepts into a structured cellular-network analysis project. Connected protocol and architecture fundamentals with practical interpretation of end-to-end telecommunications behavior.\par
{\small\textit{Tools \& Methods:} 4G Networks, LTE, RAN, EPC Concepts, Mobility, Radio Resource Concepts, Cellular Networks, Telecommunications\par}
\vspace{1pt}
\entry{5G + Edge + IoT + AI Architecture Analysis}{Certification Project \textbar{} 2026}
Evaluated how 5G connectivity, edge computing, IoT, and AI interact in service architectures and enterprise deployment scenarios. Analyzed use-case requirements, edge placement, IoT data flows, AI processing layers, and business/operational tradeoffs for connected intelligent systems.\par
{\small\textit{Tools \& Methods:} 5G, Edge Computing, IoT, Artificial Intelligence, Architecture Analysis, Use-Case Evaluation, Data Flows, Deployment Tradeoffs\par}
\vspace{1pt}
\Needspace{2\baselineskip}
\subsection{Publications}
\entry{\href{https://doi.org/10.3390/s25247651}{On Multi-Parameter Optimization and Proactive Reliability in 5G and Beyond Cellular Networks}}{2025}
\textit{Sensors 2025}\par
\entry{\href{https://doi.org/10.1109/PIMRC59610.2024.10817403}{An AI-Driven Framework for Enhancing Resilience in Propagation Models to Enable Digital Twin}}{2024}
\textit{IEEE PIMRC}\par
\entry{Data-Driven Intelligent Outage Management for High Shadowing Environments in 5G\&B Networks}{2024}
\textit{IEEE SmartNets}\par
\entry{\href{https://arxiv.org/abs/2203.15899}{Toward a Hybrid RF/Optical Lunar Communication System (LunarComm)}}{2022}
\textit{IEEE Network}\par
\entry{Towards Positioning Error Impact Characterization and Minimization in User-Centric RAN}{2022}
\textit{IEEE WCNC}\par
\entry{\href{https://doi.org/10.1016/j.aeue.2020.153244}{Unification of RF Energy Harvesting Schemes under Mixed Rayleigh-Rician Fading Channels}}{2020}
\textit{AEÜ - International Journal of Electronics and Communications}\par
\Needspace{3\baselineskip}
\section{Certifications \& Professional Development}
\Needspace{3\baselineskip}\noindent\textbf{1. AI Foundations \& Industry Overview for Telecommunication} {\small --- \textcolor{blue}{\underline{\href{https://www.coursera.org/account/accomplishments/verify/3900EVOH4I1Y}{Coursera}}} \textbar{} Completed \textbar{} June 19, 2026 \textbar{} \textcolor{blue}{\underline{\href{https://waseemraza844.github.io/portfolio/certificates/wireless/c1-ai-found-indus-overview.pdf}{Certificate}}}}\par
\textbf{Focus:} Introduces AI and machine-learning concepts in the telecommunications domain, connecting network data with automation, optimization, predictive analytics, and operational decision support. Reviews how AI can support network planning, service assurance, and customer-facing workflows.\par
{\small\textbf{Tools \& Skills:} telecom AI/ML use cases, network analytics, automation, optimization, predictive modeling concepts, operational intelligence\par}
\vspace{2pt}
\Needspace{3\baselineskip}\noindent\textbf{2. Network and Security Optimization in Telecommunication} {\small --- \textcolor{blue}{\underline{\href{https://www.coursera.org/account/accomplishments/verify/4TQ51W1OQWRE}{Coursera}}} \textbar{} Completed \textbar{} June 20, 2026 \textbar{} \textcolor{blue}{\underline{\href{https://waseemraza844.github.io/portfolio/certificates/wireless/c2-net-sec-optim-tel.pdf}{Certificate}}}}\par
\textbf{Focus:} Covers telecommunications network optimization together with security considerations that affect service performance and reliability. Emphasizes performance analysis, operational improvement, risk awareness, and protecting network services while maintaining efficient operation.\par
{\small\textbf{Tools \& Skills:} telecom network optimization, network performance, security concepts, service reliability, operational analysis, risk awareness, network improvement\par}
\vspace{2pt}
\Needspace{3\baselineskip}\noindent\textbf{3. Customer Experience in Telecommunication} {\small --- \textcolor{blue}{\underline{\href{https://www.coursera.org/account/accomplishments/verify/NXOIIASQQXKX}{Coursera}}} \textbar{} Completed \textbar{} June 28, 2026 \textbar{} \textcolor{blue}{\underline{\href{https://waseemraza844.github.io/portfolio/certificates/wireless/c3-cust-exp-telecom.pdf}{Certificate}}}}\par
\textbf{Focus:} Covers customer-experience management in telecommunications by linking service quality, customer journeys, operational performance, and analytics to satisfaction and retention. Emphasizes identifying experience gaps through telecom service data and performance indicators.\par
{\small\textbf{Tools \& Skills:} customer experience, quality-of-experience concepts, customer journey analysis, telecom analytics, churn and retention concepts, service improvement\par}
\vspace{2pt}
\Needspace{3\baselineskip}\noindent\textbf{4. 4G Network Fundamentals} {\small --- \textcolor{blue}{\underline{\href{https://www.coursera.org/account/accomplishments/verify/DLW1YN474VZJ}{Institut Mines-Télécom / Coursera}}} \textbar{} Completed \textbar{} August 2, 2026 \textbar{} \textcolor{blue}{\underline{\href{https://waseemraza844.github.io/portfolio/certificates/wireless/4g-network-fundamentals.pdf}{Certificate}}}}\par
\textbf{Focus:} Covers LTE/4G architecture from radio access through core-network functions, including cellular-network evolution, mobility, resource management, and end-to-end service operation. Emphasizes how major network components interact to deliver reliable broadband connectivity.\par
{\small\textbf{Tools \& Skills:} LTE/4G architecture, RAN and EPC concepts, mobility and handover, radio resource management, cellular networking, network-service fundamentals\par}
\vspace{2pt}
\Needspace{3\baselineskip}\noindent\textbf{5. Business Considerations for 5G with Edge, IoT, and AI} {\small --- \textcolor{blue}{\underline{\href{https://www.coursera.org/account/accomplishments/verify/SAUEVQXV5W71}{EDUCBA / Coursera}}} \textbar{} Completed \textbar{} February 18, 2026 \textbar{} \textcolor{blue}{\underline{\href{https://waseemraza844.github.io/portfolio/certificates/wireless/bus-cons-5g.pdf}{Certificate}}}}\par
\textbf{Focus:} Examines how 5G capabilities combine with edge computing, IoT, and AI to enable new enterprise and telecommunications services. Emphasizes use-case evaluation, business value, deployment considerations, and technology tradeoffs across connected intelligent systems.\par
{\small\textbf{Tools \& Skills:} 5G business cases, edge computing, IoT, AI integration, use-case analysis, deployment tradeoffs, telecom strategy\par}
\vspace{2pt}
\Needspace{3\baselineskip}\noindent\textbf{6. Machine Learning} {\small --- \textcolor{blue}{\underline{\href{https://www.coursera.org/account/accomplishments/verify/VU3Y5TBLVLZT}{Stanford Online / Coursera}}} \textbar{} Completed \textbar{} March 29, 2020}\par
\textbf{Focus:} Covers foundational machine-learning methods for supervised and unsupervised problems, including regression, classification, model evaluation, regularization, clustering, and practical model-selection considerations. Emphasizes translating data into predictive models while controlling generalization error.\par
{\small\textbf{Tools \& Skills:} supervised learning, regression, classification, regularization, clustering, model evaluation, bias-variance tradeoffs, machine-learning workflows\par}
\vspace{2pt}
\end{document}
